\documentclass[10pt,a4paper]{amsart}
\usepackage{amsmath,amssymb,amsthm,graphicx}
\usepackage{verbatim}
\newenvironment{reference}{\comment}{\endcomment}
%\newenvironment{reference}{}{}
\newenvironment{slogan}{\comment}{\endcomment}
\newenvironment{history}{\comment}{\endcomment}

% For commutative diagrams we use Xy-pic
\usepackage[all]{xy}

% We use 2cell for 2-commutative diagrams.
\xyoption{2cell}
% Theorem environments.
%
\theoremstyle{plain}
\newtheorem{theorem}[subsection]{Theorem}
\newtheorem{proposition}[subsection]{Proposition}
\newtheorem{lemma}[subsection]{Lemma}

\theoremstyle{definition}
\newtheorem{definition}[subsection]{Definition}
\newtheorem{example}[subsection]{Example}
\newtheorem{exercise}[subsection]{Exercise}
\newtheorem{situation}[subsection]{Situation}

\theoremstyle{remark}
\newtheorem{remark}[subsection]{Remark}
\newtheorem{remarks}[subsection]{Remarks}

\numberwithin{equation}{subsection}

% Macros
%
\def\lim{\mathop{\mathrm{lim}}\nolimits}
\def\colim{\mathop{\mathrm{colim}}\nolimits}
\def\Spec{\mathop{\mathrm{Spec}}}
\def\Hom{\mathop{\mathrm{Hom}}\nolimits}
\def\Ext{\mathop{\mathrm{Ext}}\nolimits}
\def\SheafHom{\mathop{\mathcal{H}\!\mathit{om}}\nolimits}
\def\SheafExt{\mathop{\mathcal{E}\!\mathit{xt}}\nolimits}
\def\Sch{\mathit{Sch}}
\def\Mor{\mathop{\mathrm{Mor}}\nolimits}
\def\Ob{\mathop{\mathrm{Ob}}\nolimits}
\def\Sh{\mathop{\mathit{Sh}}\nolimits}
\def\NL{\mathop{N\!L}\nolimits}
\def\CH{\mathop{\mathrm{CH}}\nolimits}
\def\proetale{{pro\text{-}\acute{e}tale}}
\def\etale{{\acute{e}tale}}
\def\QCoh{\mathit{QCoh}}
\def\Ker{\mathop{\mathrm{Ker}}}
\def\Im{\mathop{\mathrm{Im}}}
\def\Coker{\mathop{\mathrm{Coker}}}
\def\Coim{\mathop{\mathrm{Coim}}}

% Boxtimes
%
\DeclareMathSymbol{\boxtimes}{\mathbin}{AMSa}{"02}

%
% Macros for moduli stacks/spaces
%
\def\QCohstack{\mathcal{QC}\!\mathit{oh}}
\def\Cohstack{\mathcal{C}\!\mathit{oh}}
\def\Spacesstack{\mathcal{S}\!\mathit{paces}}
\def\Quotfunctor{\mathrm{Quot}}
\def\Hilbfunctor{\mathrm{Hilb}}
\def\Curvesstack{\mathcal{C}\!\mathit{urves}}
\def\Polarizedstack{\mathcal{P}\!\mathit{olarized}}
\def\Complexesstack{\mathcal{C}\!\mathit{omplexes}}
% \Pic is the operator that assigns to X its picard group, usage \Pic(X)
% \Picardstack_{X/B} denotes the Picard stack of X over B
% \Picardfunctor_{X/B} denotes the Picard functor of X over B
\def\Pic{\mathop{\mathrm{Pic}}\nolimits}
\def\Picardstack{\mathcal{P}\!\mathit{ic}}
\def\Picardfunctor{\mathrm{Pic}}
\def\Deformationcategory{\mathcal{D}\!\mathit{ef}}
\title{The internal hom of presheaves of modules commutes with limits and the internal hom of sheaves does too}
\author{The Stacks Project Authors}
\date{stacks-project:sites-modules.tex:local-snapshot-20261012; GNU FDL 1.2 or later}
\begin{document}
\maketitle
\noindent\textit{Editorial note.} All delivered mathematical text is the author's own, copied verbatim from the pinned Stacks Project chapter file sites-modules.tex; no statement or proof step was rewritten, reordered or removed.\par\medskip
\noindent\textit{Editorial note.} The counted span cites no labelled external item (the adjointness used is discussed in the same section), so no navigation stubs were needed.\par\medskip
\noindent\textit{Allowed standard prerequisites (editorial).}\begin{itemize}
\item Restriction to the localized category is both a left and a right adjoint. Stacks Project, Modules on Sites, section on internal hom (adjointness of j\_U\textasciicircum{}*). Gives preservation of limits by restriction, which the sectionwise computation needs.
\end{itemize}\medskip
\begin{lemma}
\label{lemma-internal-hom-commute-limits}
Internal hom and (co)limits.
Let $\mathcal{C}$ be a category and let $\mathcal{O}$ be a presheaf of rings.
\begin{enumerate}
\item For any presheaf of $\mathcal{O}$-modules $\mathcal{F}$ the functor
$$
\textit{PMod}(\mathcal{O}) \longrightarrow \textit{PMod}(\mathcal{O})
, \quad
\mathcal{G} \longmapsto \SheafHom_\mathcal{O}(\mathcal{F}, \mathcal{G})
$$
commutes with arbitrary limits.
\item For any presheaf of $\mathcal{O}$-modules $\mathcal{G}$ the functor
$$
\textit{PMod}(\mathcal{O}) \longrightarrow \textit{PMod}(\mathcal{O})^{opp}
, \quad
\mathcal{F} \longmapsto \SheafHom_\mathcal{O}(\mathcal{F}, \mathcal{G})
$$
commutes with arbitrary colimits, in a formula
$$
\SheafHom_\mathcal{O}(\colim_i \mathcal{F}_i, \mathcal{G})
=
\lim_i \SheafHom_\mathcal{O}(\mathcal{F}_i, \mathcal{G}).
$$
\end{enumerate}
Suppose that $\mathcal{C}$ is a site, and $\mathcal{O}$ is a sheaf of rings.
\begin{enumerate}
\item[(3)] For any sheaf of $\mathcal{O}$-modules $\mathcal{F}$ the functor
$$
\textit{Mod}(\mathcal{O}) \longrightarrow \textit{Mod}(\mathcal{O})
, \quad
\mathcal{G} \longmapsto \SheafHom_\mathcal{O}(\mathcal{F}, \mathcal{G})
$$
commutes with arbitrary limits.
\item[(4)] For any sheaf of $\mathcal{O}$-modules $\mathcal{G}$ the functor
$$
\textit{Mod}(\mathcal{O}) \longrightarrow \textit{Mod}(\mathcal{O})^{opp}
, \quad
\mathcal{F} \longmapsto \SheafHom_\mathcal{O}(\mathcal{F}, \mathcal{G})
$$
commutes with arbitrary colimits, in a formula
$$
\SheafHom_\mathcal{O}(\colim_i \mathcal{F}_i, \mathcal{G})
=
\lim_i \SheafHom_\mathcal{O}(\mathcal{F}_i, \mathcal{G}).
$$
\end{enumerate}
\end{lemma}

\begin{proof}
Let $\mathcal{I} \to \textit{PMod}(\mathcal{O})$, $i \mapsto \mathcal{G}_i$
be a diagram. Let $U$ be an object of the category $\mathcal{C}$.
As $j_U^*$ is both a left and a right adjoint we see that
$\lim_i j_U^*\mathcal{G}_i = j_U^* \lim_i \mathcal{G}_i$.
Hence we have
\begin{align*}
\SheafHom_\mathcal{O}(\mathcal{F}, \lim_i \mathcal{G}_i)(U)
& =
\Hom_{\mathcal{O}_U}(\mathcal{F}|_U, \lim_i \mathcal{G}_i|_U) \\
& =
\lim_i \Hom_{\mathcal{O}_U}(\mathcal{F}|_U, \mathcal{G}_i|_U) \\
& = \lim_i \SheafHom_\mathcal{O}(\mathcal{F}, \mathcal{G}_i)(U)
\end{align*}
by definition of a limit. This proves (1). Part (2) is proved in exactly the
same way. Part (3) follows from (1) because the limit of a diagram of sheaves
is the same as the limit in the category of presheaves.
Finally, (4) follow because, in the formula we have
$$
\Mor_{\textit{Mod}(\mathcal{O})}(
\colim_i \mathcal{F}_i, \mathcal{G})
=
\Mor_{\textit{PMod}(\mathcal{O})}(
\colim^{PSh}_i \mathcal{F}_i, \mathcal{G})
$$
as the colimit $\colim_i \mathcal{F}_i$ is the sheafification of
the colimit $\colim^{PSh}_i \mathcal{F}_i$ in
$\textit{PMod}(\mathcal{O})$. Hence (4) follows from (2)
(by the remark on limits above again).
\end{proof}

\end{document}
